% =============================================================================
% Chapter 2: Requirements Analysis  (10–14 pages target)
% =============================================================================
\chapter{Requirements Analysis}
\label{chap:requirements}

\section*{Introduction}

This chapter presents a thorough analysis of the requirements that guided the design and implementation of \eduspace{}.
It begins by identifying the actors who interact with the system and defining the functional and non-functional requirements.
It then presents the UML modeling artifacts---use case diagrams, use case descriptions, and system sequence diagrams---that formalize the system's expected behavior.
The chapter concludes with the task distribution and project planning.

% =============================================================================
\section{Identification of Actors and Requirements}
\label{sec:actors-requirements}

% ─────────────────────────────────────────────────────────────────────────────
\subsection{Persona Sheets (System Actors)}
\label{sec:personas}

The \eduspace{} platform defines four actors, distinguished by their authentication state and their role within a given classroom.

\begin{enumerate}[label=\textbf{A\arabic*.}]
  \item \textbf{Visitor.}
        An unauthenticated user who can access public pages (landing page, login, registration, password reset).
        A visitor becomes a \emph{User} after registering and verifying their email address.

  \item \textbf{User (Authenticated).}
        A registered and verified individual who can access the dashboard, create or join classrooms, manage their profile, and view notifications and the calendar.
        Upon entering a classroom, the user assumes one of the two classroom-level roles below.

  \item \textbf{Member.}
        A user who has joined a classroom with the \texttt{MEMBER} role.
        Members can view the stream, filter and search posts, comment, submit assignments (in Teaching classrooms), view their own grades, and participate in the group chat.

  \item \textbf{Admin.}
        A user who holds the \texttt{ADMIN} role in a classroom (the creator is automatically an Admin; other members can be promoted).
        In addition to all Member capabilities, Admins can create, edit, and delete posts, manage chapters, manage members (promote or remove), grade submissions, configure classroom settings, and enable or disable the chat.
\end{enumerate}

% ─────────────────────────────────────────────────────────────────────────────
\subsection{Functional Requirements}
\label{sec:functional-requirements}

The functional requirements of \eduspace{} are organized below by domain.
Each requirement is derived from the project specification and describes a concrete capability that the system provides to its users.

% ── Authentication & User Management ────────────────────────────────────────
\subsubsection{Authentication and User Management}

\begin{enumerate}[label=\textbf{FR-\arabic*.}, leftmargin=2.5cm]
  \item \textbf{Registration.}
        A visitor can create an account by providing an email address, a password, a first name, and a last name.
        The password is hashed using bcrypt before storage.

  \item \textbf{Email verification.}
        Upon registration, the system sends a six-digit verification code to the provided email address.
        The user must enter this code to activate the account.

  \item \textbf{Login.}
        A registered and verified user can log in with email and password.

  \item \textbf{Password reset.}
        A user who has forgotten their password can request a reset link via email.

  \item \textbf{Profile management.}
        An authenticated user can view and update their profile information.

  \item \textbf{Account deletion.}
        An authenticated user can permanently delete their account.
\end{enumerate}

\begin{figure}[H]
  \centering
  
\includegraphics[width=0.85\textwidth]{figma/01 login.png}
  \caption{Figma mockup of the login page.}
  \label{fig:figma-login}
\end{figure}

% ── Classroom Management ────────────────────────────────────────────────────
\subsubsection{Classroom Management}

\begin{enumerate}[label=\textbf{FR-\arabic*.}, leftmargin=2.5cm, resume]
  \item \textbf{Classroom creation.}
        An authenticated user can create a classroom (\emph{Teaching} or \emph{Friendly}).
        The system generates a unique six-character alphanumeric class code.

  \item \textbf{Dual classroom model.}
        The platform supports two classroom types:
        \begin{itemize}[nosep]
          \item \emph{Teaching} --- the creator and promoted users are displayed as ``Teachers,'' joiners as ``Students.'' Grading features are fully enabled.
          \item \emph{Friendly} --- the creator and promoted users are displayed as ``Admins,'' joiners as ``Members.'' Grading features are hidden entirely; the API rejects any grading-related requests.
        \end{itemize}
        Both types use the same two internal roles: \texttt{ADMIN} and \texttt{MEMBER}.

  \item \textbf{Joining a classroom.}
        Any authenticated user can join a classroom by entering its six-character class code.

  \item \textbf{Leaving a classroom.}
        Any member (except the creator) can leave a classroom voluntarily.

  \item \textbf{Classroom settings.}
        Administrators can update the classroom name, description, section, subject, cover image, and the chat toggle (enable or disable group chat).

  \item \textbf{Classroom deletion.}
        Only the original creator of a classroom can delete it.

  \item \textbf{Listing classrooms.}
        An authenticated user can view a dashboard showing all classrooms they belong to.
\end{enumerate}

\begin{figure}[H]
  \centering
  \includegraphics[width=0.85\textwidth]{figma/05 dashboard-with-classrooms.png}
  \caption{Figma mockup of the dashboard with classroom cards.}
  \label{fig:figma-dashboard}
\end{figure}

% ── Member Management ───────────────────────────────────────────────────────
\subsubsection{Member Management}

\begin{enumerate}[label=\textbf{FR-\arabic*.}, leftmargin=2.5cm, resume]
  \item \textbf{Member listing.}
        All classroom members can view the list of members, grouped by role.
        The header labels adapt to the classroom type (``Teachers \& Students'' vs.\ ``Admins \& Members'').

  \item \textbf{Role promotion.}
        An administrator can promote a \texttt{MEMBER} to \texttt{ADMIN}.

  \item \textbf{Member removal.}
        An administrator can remove any member from the classroom, except the classroom creator.
\end{enumerate}


% ── Chapter Management ──────────────────────────────────────────────────────
\subsubsection{Chapter Management}

\begin{enumerate}[label=\textbf{FR-\arabic*.}, leftmargin=2.5cm, resume]
  \item \textbf{Chapter creation.}
        Administrators can create new chapters within a classroom to organize content thematically.

  \item \textbf{Chapter reordering.}
        Administrators can reorder chapters by drag-and-drop.

  \item \textbf{Chapter editing and deletion.}
        Administrators can rename or delete chapters.
        The ``General'' chapter, which is auto-created with each classroom, cannot be renamed or deleted.
\end{enumerate}

% ── Post Management ─────────────────────────────────────────────────────────
\subsubsection{Post Management}

\begin{enumerate}[label=\textbf{FR-\arabic*.}, leftmargin=2.5cm, resume]
  \item \textbf{Post creation.}
        Administrators can create posts within a classroom.
        Each post requires a title, a type, and a chapter assignment.
        The five post types are:
        \begin{itemize}[nosep]
          \item \textbf{Study Material} --- further categorized into sub-types: Cours (lecture material), TD (tutorial exercises), TP (lab work), and Summary.
          \item \textbf{Quiz} --- a quiz or test for the class.
          \item \textbf{Question} --- a question posted to the class.
          \item \textbf{Assignment} --- includes a due date and maximum points (Teaching classrooms only).
          \item \textbf{Announcement} --- a general announcement.
        \end{itemize}
        All post types support optional text content and file attachments.

  \item \textbf{Post viewing.}
        All members can view posts.
        The \emph{Stream} tab displays posts in reverse chronological order.

  \item \textbf{Stream filtering and sorting.}
        The stream provides a toolbar with:
        \begin{itemize}[nosep]
          \item A type filter (all types, or a specific type; selecting Study Material reveals a secondary sub-type filter).
          \item A chapter filter (all chapters, or a specific chapter).
          \item Sort controls (newest first, oldest first, title A--Z, title Z--A).
          \item A search bar connected to the full-text search engine.
        \end{itemize}

  \item \textbf{Post editing and deletion.}
        Any administrator can update or delete any post in the classroom.

  \item \textbf{Commenting.}
        All members can post comments on any post.
        Comment authors can edit or delete their own comments.
        Administrators can delete any comment.
\end{enumerate}

\begin{figure}[H]
  \centering
  \includegraphics[width=0.85\textwidth]{figma/13 stream-tab.png}
  \caption{Figma mockup of the Stream tab with post feed and filter toolbar.}
  \label{fig:figma-stream}
\end{figure}

% ── Assignment and Grading Workflow ─────────────────────────────────────────
\subsubsection{Assignment and Grading Workflow}

These features are available exclusively in \emph{Teaching} classrooms.

\begin{enumerate}[label=\textbf{FR-\arabic*.}, leftmargin=2.5cm, resume]
  \item \textbf{Assignment submission.}
        Members can submit work for an assignment by providing optional text content and file attachments.

  \item \textbf{Submission grading.}
        Administrators can view all submissions for an assignment.

  \item \textbf{Assignments calendar.}
        A calendar view (month and week modes) displays assignment due dates for the classroom, providing a visual overview of upcoming deadlines.
\end{enumerate}


% ── Communication ───────────────────────────────────────────────────────────
\subsubsection{Real-Time Group Chat}

\begin{enumerate}[label=\textbf{FR-\arabic*.}, leftmargin=2.5cm, resume]
  \item \textbf{Automatic chat room creation.}
        When a classroom is created with chat enabled, a dedicated chat room is automatically provisioned.
        Each classroom has exactly one chat room.

  \item \textbf{Chat toggle.}
        Administrators can enable or disable the group chat via classroom settings.
        Disabling the chat hides the Chat tab from the interface but does not delete the chat history.

  \item \textbf{Real-time messaging.}
        Members can send text messages and share files in real time using WebSocket connections (Socket.IO).

  \item \textbf{Typing indicators.}
        The system broadcasts typing and stop-typing events to all connected members in the same chat room.
\end{enumerate}

\begin{figure}[H]
  \centering
  \includegraphics[width=0.85\textwidth]{figma/12 chat-page.png}
  \caption{Figma mockup of the real-time group chat interface.}
  \label{fig:figma-chat}
\end{figure}

% ── Notifications ───────────────────────────────────────────────────────────
\subsubsection{Notification System}

\begin{enumerate}[label=\textbf{FR-\arabic*.}, leftmargin=2.5cm, resume]
  \item \textbf{Event-driven notifications.}
        The system generates notifications for key events:
        \begin{itemize}[nosep]
          \item A new post is created in a classroom.
          \item An assignment is due within 24 hours.
          \item A submission has been graded.
          \item A new chat message is received (when the user is not active in the chat).
        \end{itemize}

  \item \textbf{Real-time delivery.}
        Notifications are pushed to connected clients in real time via Redis Pub/Sub.
        An unread count badge is displayed in the navigation bar.

  \item \textbf{Notification management.}
        Users can view their notifications, mark them as read (individually or all at once), and delete them (individually or all at once).
\end{enumerate}

\medskip

In total, 34 functional requirements were identified across various domains: authentication and user management, classroom management, member management, chapter management, post management, assignment and grading, file management, group chat, notifications, and search.

% ─────────────────────────────────────────────────────────────────────────────
\subsection{Non-Functional Requirements}
\label{sec:non-functional}

Beyond the functional capabilities, \eduspace{} must satisfy the following non-functional requirements to ensure a robust, secure, and maintainable system.

\begin{enumerate}[label=\textbf{NFR-\arabic*.}, leftmargin=2.8cm]
  \item \textbf{Performance.}
        \begin{itemize}[nosep]
          \item API response times must remain under 500\,ms for standard CRUD operations.
          \item Real-time chat messages must be delivered to all connected participants within 200\,ms under normal load.
          \item The search service must return results within 300\,ms for typical queries.
        \end{itemize}

  \item \textbf{Security.}
        \begin{itemize}[nosep]
          \item All passwords must be hashed using bcrypt before storage.
          \item Authentication is enforced via JWT tokens.
          \item WebSocket connections require a valid JWT token upon handshake.
        \end{itemize}

  \item \textbf{Scalability.}
        \begin{itemize}[nosep]
          \item The microservices architecture allows individual services to be scaled independently based on load.
          \item Each service owns its dedicated database.
        \end{itemize}

  \item \textbf{Availability and reliability.}
        \begin{itemize}[nosep]
          \item The system is containerized with Docker Compose, enabling reproducible deployments and rapid recovery.
          \item Service failures are isolated: if one microservice goes down, the rest continue to operate independently.
        \end{itemize}

  \item \textbf{Usability.}
        \begin{itemize}[nosep]
          \item The frontend provides a responsive, modern interface.
          \item Navigation is consistent across pages with a persistent sidebar and top navigation bar.
        \end{itemize}
\end{enumerate}

% =============================================================================
\section{Requirements Modeling}
\label{sec:modeling}

% ─────────────────────────────────────────────────────────────────────────────
\subsection{Use Case Diagrams}
\label{sec:use-case-diagrams}

\figref{fig:use-case-global} presents the global use case diagram of \eduspace{}, organized into seven functional packages: Authentication, Profile \& Account, Dashboard, Classroom Administration, Classroom Content, Assignments \& Grades, Chat, and Files.
The four actors---Visitor, User, Member, and Admin---interact with the system according to their privilege level.

\begin{figure}[H]
  \centering
  \rotatebox{90}{\includegraphics[width=0.85\textheight,height=\textwidth,keepaspectratio]{diagrams/use-case.pdf}}
  \caption{Global use case diagram of the \eduspace{} platform.}
  \label{fig:use-case-global}
\end{figure}

% ─────────────────────────────────────────────────────────────────────────────
\subsection{Use Case Descriptions}
\label{sec:use-case-descriptions}

This section describes the major use cases through nominal and alternative scenarios.

\subsubsection{UC-1: Register and Verify Email}

\begin{tabularx}{\textwidth}{@{} l X @{}}
  \toprule
  \textbf{Actor}          & Visitor \\
  \textbf{Precondition}   & The visitor has no existing account. \\
  \textbf{Postcondition}  & A verified user account exists; the user can log in. \\
  \midrule
  \multicolumn{2}{@{}l}{\textbf{Nominal scenario}} \\
  \midrule
  1 & The visitor opens the registration page. \\
  2 & The visitor fills in email, password, first name, and last name, then submits. \\
  3 & The system creates the account, hashes the password, generates a 6-digit code, and sends it by email. \\
  4 & The visitor enters the verification code. \\
  5 & The system marks the account as verified and redirects to the login page. \\
  \midrule
  \multicolumn{2}{@{}l}{\textbf{Alternative scenarios}} \\
  \midrule
  3a & Email already exists: the system returns an error. \\
  4a & Invalid code: the system displays an error; the user may request a new code. \\
  \bottomrule
\end{tabularx}

\subsubsection{UC-2: Login}

\begin{tabularx}{\textwidth}{@{} l X @{}}
  \toprule
  \textbf{Actor}          & Visitor \\
  \textbf{Precondition}   & The user has a verified account. \\
  \textbf{Postcondition}  & The user is authenticated and redirected to the dashboard. \\
  \midrule
  \multicolumn{2}{@{}l}{\textbf{Nominal scenario}} \\
  \midrule
  1 & The visitor opens the login page. \\
  2 & The visitor enters email and password, then submits. \\
  3 & The system validates credentials, generates a JWT token, and returns it with user profile data. \\
  4 & The frontend stores the token and redirects to the dashboard. \\
  \midrule
  \multicolumn{2}{@{}l}{\textbf{Alternative scenarios}} \\
  \midrule
  3a & Invalid credentials: the system returns a 401 error. \\
  3b & Email not verified: the system returns a 403 error and redirects to verification. \\
  \bottomrule
\end{tabularx}

\subsubsection{UC-3: Create a Classroom}

\begin{tabularx}{\textwidth}{@{} l X @{}}
  \toprule
  \textbf{Actor}          & User \\
  \textbf{Precondition}   & The user is authenticated. \\
  \textbf{Postcondition}  & A new classroom exists; the user is its creator (Admin). \\
  \midrule
  \multicolumn{2}{@{}l}{\textbf{Nominal scenario}} \\
  \midrule
  1 & The user clicks the create button on the dashboard. \\
  2 & The user fills in name, type (Teaching/Friendly), and optional fields (subject, section, description), then submits. \\
  3 & The system generates a unique 6-character class code, creates the classroom, adds the user as Admin, creates a ``General'' chapter, and publishes a \texttt{classroom.created} event. \\
  4 & The user is redirected to the new classroom. \\
  \bottomrule
\end{tabularx}

\subsubsection{UC-4: Create a Post}

\begin{tabularx}{\textwidth}{@{} l X @{}}
  \toprule
  \textbf{Actor}          & Admin \\
  \textbf{Precondition}   & The admin is in a classroom. \\
  \textbf{Postcondition}  & A new post appears in the stream; members are notified. \\
  \midrule
  \multicolumn{2}{@{}l}{\textbf{Nominal scenario}} \\
  \midrule
  1 & The admin opens the post creation form. \\
  2 & The admin selects a type, chapter, and fills in the title, content, and optional attachments. \\
  3 & If files are attached, they are uploaded to MinIO via the File Service. \\
  4 & The system verifies the admin's role, creates the post in a transaction, and publishes a \texttt{post.created} event. \\
  5 & The post appears in the stream; the Search Service indexes it; the Notification Service notifies members. \\
  \midrule
  \multicolumn{2}{@{}l}{\textbf{Alternative scenarios}} \\
  \midrule
  2a & Type is Assignment in a Friendly classroom: the API rejects the request with a 400 error. \\
  4a & User is not an Admin: the API returns a 403 error. \\
  \bottomrule
\end{tabularx}

\subsubsection{UC-5: Submit and Grade an Assignment}

\begin{tabularx}{\textwidth}{@{} l X @{}}
  \toprule
  \textbf{Actors}         & Member (submit), Admin (grade) \\
  \textbf{Precondition}   & An Assignment post exists in a Teaching classroom. \\
  \textbf{Postcondition}  & The submission is recorded and graded; the student is notified. \\
  \midrule
  \multicolumn{2}{@{}l}{\textbf{Nominal scenario}} \\
  \midrule
  1 & The member opens the assignment post and clicks submit. \\
  2 & The member optionally writes text and attaches files, then confirms. \\
  3 & The system uploads files, creates the submission, and displays a confirmation. \\
  4 & The admin views the list of submissions and selects one. \\
  5 & The admin enters points and optional feedback, then saves. \\
  6 & The system updates the submission, publishes \texttt{submission.graded}, and the student receives a notification. \\
  \midrule
  \multicolumn{2}{@{}l}{\textbf{Alternative scenarios}} \\
  \midrule
  5a & Points exceed the maximum: the API returns a 400 error. \\
  \bottomrule
\end{tabularx}

% ─────────────────────────────────────────────────────────────────────────────
\subsection{System Sequence Diagrams}
\label{sec:sequence-diagrams}

The following system sequence diagrams illustrate the interactions between the actors and the system for the major use cases.

\begin{figure}[H]
  \centering
  \includegraphics[width=\textwidth,height=0.85\textheight,keepaspectratio]{diagrams/seq-auth-login.png}
  \caption{System sequence diagram --- Authentication (Registration + Login).}
  \label{fig:seq-auth}
\end{figure}

\begin{figure}[H]
  \centering
  \rotatebox{90}{\includegraphics[width=0.85\textheight,height=\textwidth,keepaspectratio]{diagrams/seq-create-classroom.png}}
  \caption{System sequence diagram --- Classroom Creation.}
  \label{fig:seq-create-classroom}
\end{figure}

\begin{figure}[H]
  \centering
  \rotatebox{90}{\includegraphics[width=0.85\textheight,height=\textwidth,keepaspectratio]{diagrams/seq-create-post.png}}
  \caption{System sequence diagram --- Post Creation.}
  \label{fig:seq-create-post}
\end{figure}

% Note: Section 2.2.4 "Contraintes techniques" is facultatif — omitted per instructions.

% ─────────────────────────────────────────────────────────────────────────────
\subsection{Task Distribution and Planning}
\label{sec:planning}

As described in \secref{sec:organization}, the team adopted a feature-based collaboration model.
Rather than dividing responsibilities along technical layers (frontend vs.\ backend), each team member took full ownership of specific features end-to-end---from database schema and backend API to the corresponding frontend interface.
Shared concerns such as UML diagrams, the shared utility package, infrastructure configuration, and architectural decisions were handled collaboratively.

The project was organized into six major phases, spanning approximately four months from February to May 2026.
\tabref{tab:phases} summarizes the phases and their scope.

\begin{table}[H]
  \centering
  \caption{Project phases and scope.}
  \label{tab:phases}
  \small
  \begin{tabularx}{\textwidth}{@{} c l X @{}}
    \toprule
    \textbf{Phase} & \textbf{Period} & \textbf{Scope} \\
    \midrule
    1 & Feb 1--14   & Requirements analysis, UML modeling, technology selection \\
    2 & Feb 15--28  & Infrastructure setup (Docker Compose, databases, Redis, RabbitMQ, MinIO, NGINX, Elasticsearch), shared package, API Gateway \\
    3 & Mar 1--Mar 15 & Core backend services: User, Class, Content, File \\
    4 & Mar 16--Mar 31   & Communication, Notification, and Search services; inter-service event wiring (RabbitMQ) \\
    5 & Apr 1--Apr 30 & Frontend development: authentication, dashboard, classroom pages, chat, grades, calendar, settings \\
    6 & May 1--May 19  & Integration testing, bug fixes, report writing \\
    \bottomrule
  \end{tabularx}
\end{table}


% =============================================================================
\section*{Conclusion}

This chapter defined the scope of \eduspace{} through a systematic requirements analysis.
A total of 34 functional requirements were identified across varoious domains, complemented by five non-functional requirements covering performance, security, scalability, availability, maintainability, and usability.
The UML modeling artifacts formalize the system's expected behavior and interactions.
The task distribution demonstrate how the two-person team organized the work across a five-month timeline.

The next chapter presents the system design, beginning with the architectural paradigm and proceeding to the object-oriented design artifacts.
